\documentclass[10pt]{article}
% Packages
\usepackage{makecell} % For thicker lines
\usepackage{setspace}  % Controls line spacing
\usepackage{hhline}   % For double horizontal lines
\usepackage{colortbl}  % Add colour to LaTeX tables
\usepackage[T1]{fontenc}  % Choice of font encodings
\usepackage{tgtermes}  % Loads the TeX Gyre Termes font
\usepackage{siunitx}  % A comprehensive (SI) units package
\usepackage{tabularx, booktabs} % For advanced table layout
\usepackage{url}  % Verbatim with URL-sensitive line breaks
\usepackage{authblk}  % For author and affiliation management
\usepackage{natbib}  % A package for bibliographies and citations
\usepackage{graphicx}  % Enhances LaTeX's built-in graphics features
\usepackage{listings}  % Typeset programming code within the document
\usepackage{amssymb}  % Mathematical symbols
\usepackage[nottoc]{tocbibind}  % Adds entries to the table of contents
\usepackage{xcolor}  % Provides easy driver-independent access to colors
\usepackage{microtype}  % Improves the spacing between words and letters
\usepackage{enumitem}  % Control layout of itemize, enumerate, description
\usepackage{tocloft}  % Controls the typographic design of table of contents, etc.
\usepackage[breaklinks,linktocpage]{hyperref}  % Creates hyperlinks in your document
\usepackage[font=small,skip=7pt,labelfont=bf]{caption}  % Customising captions in floating envs

% Adjust the page margins in the bibliography
\let\oldthebibliography=\thebibliography
\let\endoldthebibliography=\endthebibliography
\renewenvironment{thebibliography}[1]{%
  \begin{oldthebibliography}{#1}%
    \raggedright%
    }{%
  \end{oldthebibliography}%
}

\setlength\bibindent{0pt}

% Optional options
% \usepackage{background} % Creates a DRAFT background image on all pages
% \backgroundsetup{contents=Preprint, opacity=0.25, color=gray} % Adds a watermark to the document
% This command changes the line spacing to double.
% ? Needed for reviews/drafts
% \doublespacing

% Custom colours
\definecolor{codegreen}{rgb}{0,0.5,0}
\definecolor{codegray}{rgb}{0.4,0.4,0.4}
\definecolor{codepurple}{rgb}{0.58,0,0.82}
\definecolor{backcolour}{rgb}{0.96,0.96,0.96}
\definecolor{lightgray}{gray}{0.8}

\lstdefinelanguage{JavaScript}{
  keywords={break, case, catch, continue, debugger, default, delete, do, else, finally, for, function, if, in, instanceof, new, return, switch, this, throw, try, typeof, var, void, while, with},
  morecomment=[l]{//},
  morecomment=[s]{/*}{*/},
  morestring=[b]',
  morestring=[b]",
  sensitive=true
}

% Listing styles
\lstdefinestyle{mystyle}{
  frame=tb,
  tabsize=2,
  captionpos=b,
  numbers=left,
  framerule=0pt,
  numbersep=5pt,
  showtabs=false,
  breaklines=true,
  keepspaces=true,
  showspaces=false,
  framextopmargin=6pt,
  framexbottommargin=6pt,
  showstringspaces=false,
  breakatwhitespace=false,
  keywordstyle=\color{purple},
  commentstyle=\color{codegreen},
  stringstyle=\color{codepurple},
  numberstyle=\tiny\color{codegray},
  basicstyle=\ttfamily\footnotesize,
  backgroundcolor=\color{backcolour}}
\lstset{style=mystyle}

% ! Custom template commands
% Add a vertical space after section numbers in ToC
\renewcommand\cftsecafterpnum{\vskip8pt}
% Changes the title of the list of listings
\renewcommand{\lstlistlistingname}{List of \lstlistingname s}
% Changes the title of the bibliography
\renewcommand{\bibname}{References}
% Changes the title of the table of contents
\renewcommand{\contentsname}{Table of Contents}
% Changes the leader between section and page numbers in ToC
\renewcommand{\cftsecleader}{\cftdotfill{\cftdotsep}}
\newcommand{\floatcaption}[2]{\caption[#1.]{#1~#2.}}

% Custom template settings
\captionsetup{justification=centering}  % All captions will be centered
\hypersetup{
  colorlinks = true,
  urlcolor = blue,
  linkcolor = blue,
  citecolor = blue,
  breaklinks = true
}
\PassOptionsToPackage{hyphens}{url}
\urlstyle{same}
\def\Urlmuskip{0mu plus 1mu}
\def\UrlBreaks{\do\/\do-}
\def\UrlBigBreaks{\do\/\do-\do:\do.}
\setlist[itemize]{noitemsep, topsep=0pt, partopsep=0pt}



\begin{document}
\pagenumbering{roman}
\counterwithin{lstlisting}{section}
\counterwithin{figure}{section}
\counterwithin{table}{section}
\setlength{\footskip}{65pt}

% ! ===========================
% ! # MARK: Title, author, etc.
% ! ===========================

\title{\textbf{Death is an Engineering Challenge}}
\author[1]{Daniel Burger}
\affil[1]{\textbf{Eightsix Science Ltd}}
\affil[ ]{\href{mailto:daniel@eightsix.science}{daniel@eightsix.science}}
\author[2]{Gabriel Cunha}
\affil[2]{\textbf{Tufts University}}
\affil[ ]{\href{mailto:gabriel.cunha@tufts.edu}{gabriel.cunha@tufts.edu}}
\author[3]{Izumi Handa}
\affil[3]{\textbf{Panda Lab Inc}}
\affil[ ]{\href{mailto:izumi.handa@pandalab.jp}{izumi.handa@pandalab.jp}}
\author[4]{Masataka Watanabe}
\affil[4]{\textbf{University of Tokyo}}
\affil[ ]{\href{mailto:watanabe@sys.t.u-tokyo.ac.jp}{watanabe@sys.t.u-tokyo.ac.jp}}

\date{\textit{\today}}
\maketitle
\thispagestyle{empty}

\begin{sloppypar}

  \begin{figure}[ht]
    \centering
    \includegraphics[width=\textwidth]{figures/cover.png}
    \label{fig:cover}
    \caption{\% TODO: The four-dimensional process world line of a conscious substrate going through space-time.}
  \end{figure}
  \newpage

  % ! ================
  % ! # MARK: Abstract
  % ! ================


  \begin{abstract}
    We view death as the irreversible destruction of consciousness’s physical and dynamic processes and frame it as a manageable systems problem solvable through engineering. By that, we propose synconetics, a new discipline dedicated to sustaining conscious continuity with current and near-term technologies that we can empirically test. This essay outlines the principles of synconetics and introduces two practical approaches potentially capable of achieving this goal within the next twenty years.
  \end{abstract}

  \pagebreak
  \pagenumbering{Roman}
  \tableofcontents
  \pagebreak
  \listoffigures
  \pagebreak
  % \listoftables
  % \pagebreak
  % \addcontentsline{toc}{section}{\lstlistlistingname}
  % \lstlistoflistings
  % \pagebreak
  \pagenumbering{arabic}

  % ! ====================
  % ! # MARK: Introduction
  % ! ====================

  \section{Introduction}
  \label{sec:introduction}

  Months of collaborative alignment within our working group, involving rigorous debate and scrutiny of prior schools of thought, have shown fundamental shortcomings in conventional approaches to solving death. As a result of this critical process, we introduce synconetics. This essay presents the first formal articulation of this new discipline, outlining a cohesive paradigm born from our shared analysis and convictions. However, we stress that this represents an initial definition of an emerging field; the written work presented herein, as will become apparent throughout the essay, requires significant refinement through broader engagement and critique. Recognising this necessity, we invite collaborators from academia and industry to join us in challenging, developing, and advancing synconetics.

  At its core, synconetics is a scientific discipline focused on understanding the fundamental aspects of death and overcoming it through the development of synthetic consciousness mechanics. This mechanics provides a framework for engineering interventions designed to sustain the physical-dynamical processes of consciousness using resilient substrates. By grounding its approach in empirically testable systems and near-term technologies over speculative and non-falsifiable approaches, synconetics aims to bypass philosophical speculation about the true nature of consciousness and its cessation (i.e., death). It establishes consciousness continuity preservation as a concrete engineering challenge, forming the discipline’s foundational principle.

  \subsection{First Principles of Death}
  \label{sec:first-principles}

  Death is not an inevitability but a contingent failure of the organised physical system that instantiates consciousness (i.e., the conscious substrate). In principle, physical processes are manipulable: no fundamental law prohibits the indefinite persistence of self-maintaining systems under engineered conditions. Therefore, we define dying operationally as the irreversible destruction of the processes sustaining conscious continuity over time.

  \begin{figure}[ht!]
    \centering
    \includegraphics[width=\textwidth]{figures/4D-process-world-line-ending.png}
    \caption{\% TODO: The four-dimensional process world line of a conscious substrate, what it means if it decays (i.e., death, irreversible loss of conscious continuity, the past is immutable), and what it means when we preserve it over another substrate through substrate independence.}
    \label{fig:process-world-line-ending}
  \end{figure}

  The term conscious substrate denotes any physical medium whose organisation and dynamics have been shown to support conscious experience. Biological substrates, such as the central nervous system (CNS), remain the only empirically confirmed examples. This definition does not restrict consciousness to biology but acknowledges current empirical constraints. Individuals possess direct, if subjective, evidence of their own conscious continuity (e.g., “I think, therefore I am”), providing a provisional anchor for engineering objectives.

  Death marks the irreversible termination of the four-dimensional process-world-line: the spatiotemporal trajectory of physical states within the conscious substrate that collectively underpins conscious continuity (see \autoref{fig:process-world-line-ending} for a schematic visualisation). Irreversibility reflects thermodynamic reality; an increase in entropy erodes the recoverability of prior states. Synconetics addresses this through open-system engineering, countering entropy, for example, via energy/matter exchange and error correction. While the exact organisational level essential for consciousness remains unclear (whether cellular, molecular, atomic, etc.), engineering pragmatism prevails: we do not need to fully understand a system to work productively with it and utilise what we have at hand.

  \subsection{Critiquing Conventional Paradigms}
  \label{sec:new-paradigm}

  \begin{figure}[ht!]
    \centering
    \includegraphics[width=\textwidth]{figures/other-approaches.png}
    \caption{\% TODO: Show the line of the four-dimensional process world line compared to how other approaches are doing it.}
    \label{fig:other-approaches}
  \end{figure}

  Current approaches to overcoming death, whether by delaying or reversing substrate degradation through conventional longevity research (e.g. through lifestyle interventions, genetic engineering, etc.) or by pursuing true substrate independence via computational abstraction (for instance, mind uploading or whole brain emulation), fail to address a fundamental issue: the fragility of the four-dimensional process-world-line that maintains continuous consciousness. Synconetics seeks to safeguard the physical-dynamical continuity of this process while potentially enabling measured transitions to synthetic substrates that promise greater resilience than purely biological ones.

  Mind uploading (MU) and Whole Brain Emulation (WBE) assume that consciousness can survive as a computational abstraction. This is an unverified leap that may pose an existential threat if consciousness proves inseparable from its original physical substrate. Therefore, synconetics adopts a cautious stance in that regard. Synthetic substrates may indeed be more robust, but only with thorough, incremental validation that does not break continuity. By integrating advanced elements such as brain-machine interfaces (BMIs) in ways that preserve the ongoing dynamics of the central nervous system, synconetics aims to maintain an unbroken four-dimensional process-world-line while testing the feasibility of eventually instantiating consciousness in substrates beyond standard biology. Three unresolved challenges exemplify the risks of premature reliance on MU and WBE:

  \begin{enumerate}
    \item \textbf{Philosophical Zombies}: MU and WBE’s assumptions of substrate independence, the idea consciousness can “run” on arbitrary hardware, lack empirical evidence. Computational models simulate neural correlates (e.g., firing rates, connectomes) but cannot verify whether subjective experience persists post-transfer. This leap conflates necessary conditions (physical dynamics like ion gradients and metabolic cycles) with sufficient conditions (unproven for digital substrates). Destructive uploading risks terminating the original process-world-line without verifiable assurance of qualia preservation, effectively gambling one’s existence on unconfirmed substrate equivalency. Even if functional continuity is achieved (à la Chalmers’ ‘gradual replacement’ of biological substrates with digital ones \citep{chalmers_conscious_1998}), in the absence of empirical methods to confirm subjective survival, the replica’s consciousness remains an untestable conjecture; no rigorous discipline should endorse this gamble.

    \item \textbf{Scale Separability}: Emulating brains at biologically relevant scales demands energy potentially exceeding Earth’s projected budget for centuries. Detailed simulations, such as those analysed in foundational roadmaps \citep{bostrom_whole_2008}, could model far less than 1\% of a human brain’s synaptic activity and often required slowdowns exceeding 1,000\(\times\) biological time, a trivial fraction of real-time biological efficiency (20W vs. megawatts). Crucially, the brain’s dynamics blend deterministic chaos and stochastic noise as quantum-scale fluctuations (e.g., ion channel gating) propagate into macroscopic neural activity. This “web of causality” \citep{watanabe_biological_2022}, though perhaps shorter-lived than classical models suggest, defies reduction to deterministic simulations. Worse, fidelity requirements escalate exponentially if consciousness depends on quantum effects (e.g., microtubule coherence as proposed by Orch-OR). To simulate reality faithfully at such scales, we risk asymptotically approaching the impossible task of simulating the universe within itself (i.e., infinite recursion).

    \item \textbf{Teleportation Paradox}: Most MU and WBE approaches rely on destructive scanning or non-destructive copying to transfer consciousness onto a computational substrate. Both methods sever the original’s causal continuity: destructive scans terminate the four-dimensional process-world-line outright, while non-destructive copies spawn a parallel entity divorced from your subjective stream. Even “perfect” replication creates a new four-dimensional trajectory; you do not experience the replica’s existence. Critics dismiss this as philosophical nitpicking (“I don’t care if it’s a copy, as long as it thinks it’s me.”). But survival hinges on \emph{your} conscious continuity, not a replica’s beliefs. Even if half of humanity accepted copy-based “survival,” the other half would reject existential roulette. Synconetics prioritises solutions that preserve the four-dimensional process-world-line outright, ensuring no one is forced to gamble their existence on untested metaphysics.
  \end{enumerate}

  \begin{figure}[ht!]
    \centering
    \includegraphics[width=\textwidth]{figures/4d-trajectory-disembodiment.png}
    \caption{\% TODO: Showing a graph similar to the one from the old draft in 4D, but with the disembodiment of the conscious substrate on the y-axis and the time on the x-axis. What it means to disembody it from its bodily substrate.}
    \label{fig:4d-trajectory-disembodiment}
  \end{figure}

  Critics may accuse synconetics of biological bias, but this misunderstands the burden of proof: synthetic substrates must demonstrate conscious continuity before replacing biology. Until then, privileging verified systems is engineering prudence, not chauvinism. The imperative is clear: abandon replication metaphysics and address the substrate fragility of the four-dimensional process-world-line.

  \subsection{Principles of Synconetics}
  \label{sec:principles}

  Synconetics emerges as a novel field defined by its foundational principles as shown in \autoref{tab:synconetics-principles}. We deliberately positioned it as a new, transdisciplinary scientific discipline (not confined to being a subdiscipline of any single existing field) to actively counter mother bias: the tendency for emerging disciplines to adopt the conceptual limitations and epistemic blind spots of their “parent” fields (e.g., neuroscience’s focus on biological substrates or computer science’s digital abstractionism). This positioning allows synconetics to critically integrate and innovate across multiple domains, creating a unique framework for addressing the challenges of consciousness continuity. The term synconetics refers to the core scientific discipline, while synthetic consciousness mechanics denotes the applied engineering methodologies for interfacing with, stabilising, repairing, or transitioning consciousness-supporting substrates while maintaining the four-dimensional process-world-line.

  \begin{table}[ht!]
    \centering
    \begin{tabular}{p{0.3\textwidth}p{0.65\textwidth}}
      \toprule
      \textbf{Principle}                            & \textbf{Description}                                                                                                                                                                                                               \\
      \midrule
      Process-World-Line Fidelity                   & Prioritises the preservation of the unbroken four-dimensional causal chain of the physical dynamics underlying consciousness.                                                                                                      \\
      \midrule
      Non-Destructive Transition                    & Employs gradual replacement or repair validated by first-person reportability, explicitly rejecting destructive scanning or copying methods.                                                                                       \\
      \midrule
      Substrate Agnosticism with Physical Grounding & Focuses on necessary physical dynamics rather than material composition (biological vs. synthetic), while requiring substrates to operate as thermodynamically open systems capable of sustaining consciousness-critical dynamics. \\
      \midrule
      Empirical Primacy                             & Prioritises testable interventions and falsifiable claims over philosophical speculation. Validates continuity through physical metrics and first-person reports, avoiding non-falsifiable theories of consciousness.              \\
      \midrule
      Near-Term Engineering Urgency                 & Leverages existing technologies and incremental advancements rather than awaiting speculative breakthroughs. Focuses on iterative progress using 2020s-era tools while planning for future scalability.                            \\
      \bottomrule
    \end{tabular}
    \caption{Core Principles of Synconetics}
    \label{tab:synconetics-principles}
  \end{table}

  \subsection{Nomenclature and Definitions}
  \label{sec:nomenclature}

  To ensure operational rigour and circumvent the conceptual ambiguities often found in related fields or the aforementioned conventional paradigms, synconetics employs an initial lexicon grounded directly in its foundational principles. The following core terms delineate the operational scope and key components:

  \begin{itemize}
    \item \textbf{Synthetic Consciousness Substrate (SCS)}: Any thermodynamically open physical system (biological, biohybrid or artificial) engineered or verified to be capable of sustaining the critical dynamics underlying conscious experience. Validation requires a demonstration of neural-phenomenological isomorphism with biological benchmarks. This standard reflects the profound challenge of mapping physical dynamics to subjective experience supporting millisecond-scale feedback loops essential for integrated cognitive functions, confirmed via first-person reporting during phased non-destructive integration. While achieving full isomorphism represents the ultimate validation target, practical development will involve iterative verification focusing on measurable functional equivalence, preservation of critical dynamical properties, and consistent first-person phenomenological reports.

    \item \textbf{Synthetic Consciousness Migration (SCM)}: Protocols enabling the non-destructive migration of conscious processes between substrates (e.g., from biological to SCS or between SCSs), ensuring the causal topology of the four-dimensional process-world-line remains intact. Distinct from destructive ‘scan-and-copy’ paradigms associated with conventional mind uploading, SCM focuses exclusively on methods like gradual, plasticity-aligned replacement or seamless BMI-mediated integration. The objective is to facilitate the transition to potentially more resilient or adaptable substrates while guaranteeing that the operational conditions for Continuity (see \autoref{eq:continuity}) are met throughout the process, verified by both continuous physical monitoring and first-person phenomenological reports.

    \item \textbf{Synthetic Consciousness Interfacing (SCI)}: The development and application of robust, bidirectional input/output systems connecting an SCS to external environments (physical, virtual, or hybrid) or other systems designed to operate without compromising substrate integrity or processual continuity. The primary goal is to provide the conscious entity with effective agency (the means to perceive, interact and act coherently) while also ensuring the quality and richness of interaction necessary for long-term psychological well-being and a preserved sense of self.

    \item \textbf{Processual Continuity (PC)}: Defined as the uninterrupted persistence of the causal topology (the specific, dynamic physical processes) constituting an individual’s four-dimensional process-world-line, considered essential for maintaining their unique conscious identity and subjective experience. Operationally, this continuity \( \mathcal{C} \) is maintained during an intervention or process \( \mathcal{I} \) over a time interval \( [t_0, t_1] \) if the trajectory of the system’s critical state variables \( S_{\text{crit}}(t) \) evolves within the bounds of its adaptive capacity. This condition can be formalised by requiring that the magnitude of the rate of change of these critical variables does not exceed a state-dependent threshold \( \Lambda_{\text{adapt}} \) at any point during the process:
          \begin{equation}
            \mathcal{C}(\mathcal{I}, [t_0, t_1]) \iff \forall t \in [t_0, t_1] : \left\| \frac{dS_{\text{crit}}}{dt}(t) \right\|_{\mathcal{S}} \le \Lambda_{\text{adapt}}(S_{\text{crit}}(t))
            \label{eq:continuity}
          \end{equation}
          Where:
          \begin{itemize}
            \item \( \mathcal{C}(\mathcal{I}, [t_0, t_1]) \) represents the proposition that Processual Continuity holds for the intervention \( \mathcal{I} \) over the interval \( [t_0, t_1] \).
            \item \( S_{\text{crit}}(t) \in \mathcal{S} \) is the state vector in a high-dimensional state space \( \mathcal{S} \), capturing the critical physical and dynamical variables (e.g., neural firing rates, synaptic weights, metabolic concentrations) necessary for consciousness at time \( t \). Identifying the components of \( S_{\text{crit}} \) and the appropriate state space \( \mathcal{S} \) is a major ongoing challenge for neuroscience and synconetics.
            \item \( \frac{dS_{\text{crit}}}{dt}(t) \) is the velocity vector in the state space \( \mathcal{S} \) at time \( t \), representing the instantaneous rate of change of the critical state variables. This formulation assumes sufficient smoothness for differentiability, which may be an idealisation of complex biological dynamics.
            \item \( \| \cdot \|_{\mathcal{S}} \) denotes a suitable norm defined on the tangent space of \( \mathcal{S} \), quantifying the magnitude (speed) of the state change. The choice of the norm (e.g., Euclidean \( L^2 \), maximum \( L^\infty \)) depends on which aspects of the state change are most critical to functional integrity and are subject to empirical investigation.
            \item \( \Lambda_{\text{adapt}}: \mathcal{S} \to \mathbb{R}^+ \) is a positive scalar function representing the state-dependent adaptive capacity threshold. It quantifies the maximum speed of change the system can tolerate at state \( S_{\text{crit}}(t) \) by invoking adaptive mechanisms (e.g., neural plasticity, homeostasis) without losing functional integrity or phenomenological continuity. Characterising \( \Lambda_{\text{adapt}} \) empirically for different states and interventions is a key experimental target for SCM protocols.

                  Crucially, rigorously demonstrating that this condition holds—specifically, correlating the subjective experience of continuity (verified via first-person reports) with the objective physical dynamics captured by \(S_{\text{crit}}\) and \( \Lambda_{\text{adapt}} \)—represents a core, currently unsolved challenge for synconetics. While first-person reports and functional assessments serve as essential pragmatic proxies during development, achieving the full validation demanded by our principles requires the future development of objective, real-time measures capable of tracking these critical state dynamics and their adaptive limits.
          \end{itemize}

          This inequality (\autoref{eq:continuity}) thus formalises the core engineering constraint of synconetics: the rate at which the system is forced to change by an intervention \( \mathcal{I} \) must remain below the rate at which the system can successfully adapt to that change.
  \end{itemize}

  % ! ==============================
  % ! # MARK: Methods and Approaches
  % ! ==============================

  \section{Methods and Approaches}
  \label{sec:methods}

  The aforementioned principles and definitions require translation into practical engineering applications. This section illustrates how synconetics informs concrete research and development efforts by moving beyond its theoretical frameworks. It details two distinct approaches currently being pursued by the co-authors in academia and industry. These approaches serve as initial examples of SCM in action, showcasing verifiable strategies aimed at addressing the significant challenge of preserving the four-dimensional process-world-line.

  \subsection{Approach 1: Ectopic Cognitive Preservation}
  \label{sec:daniel-approach}

  The approach termed ‘Ectopic Cognitive Preservation’ (ECP), under development by Eightsix Science Ltd (a venture co-founded by co-author Daniel Burger), exemplifies a synconetics methodology. As a practical application of SCM, it focuses on ensuring the physical continuity of the biological substrate through gradual, technologically mediated replacement, eventually aiming to create a resilient SCS. Its core technical proposal involves the progressive, piecemeal substitution of existing biological brain tissue with biohybrid neural grafts. These grafts are intended as constructs of living neural tissue, potentially derived from the patient’s own induced pluripotent stem cells (autologous iPSCs) differentiated into appropriate neural lineages to circumvent immune rejection, integrated during advanced tissue engineering with micro- or nano-scale electronic components. These integrated elements serve various functions, such as sensing local activity, providing targeted stimulation, offering structural support, or facilitating metabolic exchange during cell maturation. Achieving and rigorously verifying true functional equivalence between the original tissue and the graft (requiring not merely basic neuronal firing but complex network dynamics, synaptic plasticity profiles, and the preservation of identity-critical information patterns necessary for maintaining the unique four-dimensional process-world-line) represents a monumental yet central challenge for this approach.

  \begin{figure}[ht]
    \centering
    \includegraphics[width=\textwidth]{figures/ecp-brain-replacement.png}
    \label{fig:ecp-replacement}
    \caption{\% TODO: Visualisation of biohybrid brain replacement.}
  \end{figure}

  ECP’s commitment to continuity critically hinges on the principle of gradualism, designed to leverage the brain’s inherent plasticity (\textit{cf.} Principle 4, Physical Realisability via plasticity) and capacity for functional reorganisation, analogous to adaptations observed in response to slow-growing lesions like low-grade gliomas. The core hypothesis defines the careful management of the rate of replacement, ensuring the condition for Processual Continuity (\autoref{eq:continuity}) is met by keeping changes within the adaptive threshold \( \Lambda_{\text{adapt}} \). Neural information processing and functional roles can migrate or be re-encoded within the new substrate without disrupting the overall continuity of cognitive processes and crucial conscious experience. This reliance on plasticity, while biologically plausible, carries inherent risks regarding the fidelity of information preservation; indeed, whether plasticity can guarantee that specific memories, learned skills, and personality nuances constituting the individual’s identity are faithfully maintained (as opposed to merely enabling functional adaptation) represents a central, unproven hypothesis for ECP. This requires deep theoretical understanding and empirical validation via first-person reports alongside behavioural assessments and physical monitoring. Therefore, methodologies for precisely controlling gradual silencing and real-time monitoring of graft integration and functional takeover are critical research components.

  \begin{figure}[ht]
    \centering
    \includegraphics[width=\textwidth]{figures/ecp-vr-integration.png}
    \label{fig:ecp-vr-integration}
    \caption{\% TODO: Visualisation of VR integration.}
  \end{figure}

  The initial outcome targeted by ECP is a rejuvenated, potentially enhanced biological or biohybrid brain residing within the original cranium. This enhanced substrate, progressively composed of the new graft material integrated with embedded electronics, primarily aims to halt or reverse age-related degradation of the brain itself, thereby addressing a primary failure mode contributing to the cessation of the process-world-line. The integrated electronics could also offer inherent capabilities for advanced SCI, enabling seamless integration with virtual or augmented reality environments without requiring separate invasive procedures at a later stage.

  The ultimate, more radical goal of ECP involves the surgical explantation of this fully replaced biohybrid brain. The sustained long-term biohybrid brain, via an advanced, closed-loop whole-brain perfusion system providing a meticulously controlled physiological environment ex vivo, would be embedded within sophisticated virtual environments through high-bandwidth SCI channels derived from the integrated electronics. This step aims to achieve complete decoupling from the vulnerabilities of the original biological body, enhancing resilience. Realising stable, long-term ex vivo maintenance presents immense technical hurdles, demanding perfect replication of complex physiological conditions. Furthermore, the profound ethical and psychological implications of explantation and existence within a potentially constrained virtual reality necessitate careful consideration beyond mere technical feasibility, touching upon questions of identity, well-being and the nature of experience itself (we discuss more about this in \autoref{sec:roadmap}).

  \begin{figure}[ht]
    \centering
    \includegraphics[width=\textwidth]{figures/ecp-brain-explant.png}
    \label{fig:ecp-brain-explant}
    \caption{\% TODO: Visualisation of brain explantation.}
  \end{figure}

  Methodologically, ECP aligns directly with the foundational principles of synconetics. It confronts biological death as a contingent substrate failure problem, addressing it through a tangible engineering methodology grounded in \textbf{Empirical Primacy} (Principle 4) by integrating advanced tissue engineering, grafting, BMI, and perfusion systems based on established science rather than speculation. Its defining characteristic is the explicit commitment to \textbf{Non-Destructive Transition} (Principle 2), prioritising Processual Continuity through gradual replacement, meticulously designed to preserve the four-dimensional process-world-line and avoid the existential risks of destructive copying. This approach relies on the brain’s known plasticity to enable the gradual integration required by Principle 2, consistent with the focus on necessary physical dynamics outlined in \textbf{Substrate Agnosticism with Physical Grounding} (Principle 3). Furthermore, ECP adheres to a core synconetic tenet by squarely focusing on survival and resilience (halting degradation and enabling repair), subordinating potential enhancements to the non-negotiable goal of continuity, even if this specific constraint is not separately enumerated in the table. Its roadmap inherently generates intermediate technologies: advanced neural simulation, high-fidelity graft production, and progressive replacement techniques. These offer significant near-term therapeutic and research value, providing a pragmatic pathway for development and funding aligned with the principle of \textbf{Near-Term Engineering Urgency} (Principle 5).

  \subsection{Approach 2: Continuity via Hemispheric Integration}
  \label{sec:chi}

  A different methodology within the synconetics framework, termed Continuity via Hemispheric Integration (CHI) and conceptually developed by co-author Professor Masataka Watanabe at the University of Tokyo, leverages insights from split-brain research to achieve non-destructive substrate transition. This approach posits that consciousness, while unified, relies on the coordinated activity of both cerebral hemispheres. Clinical evidence demonstrates that separating the hemispheres can result in two distinct streams of consciousness, which can subsequently re-fuse if reconnected. CHI proposes to exploit this phenomenon by using an advanced BMI to functionally separate the biological hemispheres, linking each to a dedicated, initially non-conscious, Artificial Hemispheric Complement (AHC): a form of synthetic consciousness substrate that must be rigorously engineered and validated to support the requisite dynamics for consciousness. This strategy fundamentally differs from classical mind uploading, aiming to migrate conscious processes gradually through functional integration and redundancy rather than destructive scanning and replication.

  The proposed CHI process commences with the surgical implantation of a high-bandwidth, segmented BMI designed to functionally isolate the native hemispheres (akin to a corpus callosotomy) while enabling bidirectional communication between each biological hemisphere and its paired AHC. This initial separation is expected to induce a temporary state of dual consciousness similar to that observed in split-brain patients; this represents a significant but theoretically manageable challenge. While this induced duality presents a profound phenomenological challenge and appears, at first glance, inconsistent with maintaining a single, unbroken process-world-line as stipulated by Principle 1, the CHI hypothesis posits that redundancy across the integrated bio-artificial system is sufficient to preserve the necessary conditions for overall conscious continuity, critically preventing irreversible cessation. Subsequently, through processes analogous to neural plasticity and potentially guided by targeted stimulation protocols, functional pathways and informational representations (including memories and learned behaviours) are hypothesised to integrate between each biological hemisphere and its corresponding AHC. This ‘neural routing’ aims to establish the AHCs as functional extensions and eventual successors to their biological counterparts.

  Continuity during this critical integration phase is predicated on redundancy. As function integrates into an AHC, the combined bio-artificial system for that hemisphere maintains processing. Crucially, even if one biological hemisphere were to fail or degrade significantly, the overall conscious process is hypothesised to persist uninterrupted within the remaining biological hemisphere and its fully integrated AHC partner, mediated by the inter-hemispheric capabilities of the BMI itself. Memory transfer is envisaged via mechanisms involving both active recall (mirrored and stored in the AHC) and potentially stimulated recall (akin to Penfield’s findings), with protocols allowing for user oversight regarding sensitive memories.

  The final stage occurs once both AHCs have achieved sufficient functional integration to fully support the dynamics previously handled by their biological partners. As the biological hemispheres naturally cease function or are allowed to terminate, the individual’s conscious process-world-line persists entirely within the paired AHCs. These two AHCs are then functionally merged via the BMI or direct inter-AHC connection, restoring a unified conscious experience within a purely synthetic substrate (SCS). This final substrate, potentially operating within a secure physical location, could interact with external reality via advanced Synthetic Consciousness Interfacing (SCI), such as immersive virtual environments or robotic avatars, thereby decoupling from biological fragility.

  Methodologically, CHI aligns with the principles of synconetics. It addresses death as a substrate failure, circumventing it via engineered transition (\textbf{Substrate Agnosticism with Physical Grounding}, Principle 3, requiring AHCs to sustain necessary dynamics). It employs a tangible \textbf{engineering methodology} centred on advanced BMIs and leveraging neuroplasticity (\textbf{Empirical Primacy}, Principle 4, building on split-brain data and testable integration). Its core design is predicated on \textbf{Non-Destructive Transition} (Principle 2), aiming to preserve the four-dimensional process-world-line through redundancy and gradual integration, explicitly avoiding destructive methods. While the transition involves profound alteration, the goal is strict adherence to \textbf{Process-World-Line Fidelity} (Principle 1) by ensuring no complete cessation of conscious processing occurs. The approach leverages foreseeable advancements in BMI technology and neuroscience (\textbf{Near-Term Engineering Urgency}, Principle 5). However, CHI faces immense technical challenges, including the development of BMIs with unprecedented fidelity and bidirectional integration capabilities, the design of AHCs capable of replicating essential neural dynamics, the verification of successful integration and memory transfer, and the profound ethical and phenomenological questions surrounding the induced split-consciousness phase. Despite these hurdles, CHI represents a compelling, continuity-focused alternative within the synconetics paradigm.

  % ! =====================================
  % ! # MARK: Feasibility and Opportunities
  % ! =====================================

  \section{Feasibility and Opportunities}
  \label{sec:feasibility}

  The methodologies outlined in this essay, ECP and CHI, represent complementary pathways toward solving death. ECP adopts a biological-first strategy, incrementally transitioning to biohybrid substrates while utilising existing neurosurgical and regenerative techniques. CHI prioritises synthetic substrates from inception, using split-brain dynamics as a gateway to engineer synthetic consciousness systems. These approaches are not mutually exclusive; instead, they form a strategic duality. ECP’s gradual replacement mitigates near-term risks by preserving biological continuity, while CHI’s synthetic focus probes the boundaries of substrate independence. Critically, advances in one methodology directly inform the other. For instance, CHI’s synthetic substrate validation protocols could eventually enable ECP to transition fully to artificial systems, while ECP’s neurointegration techniques refine CHI’s bio-synthetic interfaces.

  \begin{figure}[ht!]
    \centering
    \includegraphics[width=\textwidth]{figures/complementary-pathways.png}
    \caption{\% TODO: Complementary Synconetic Strategies. ECP (top path) focuses on incremental biohybrid replacement, leveraging biological mechanisms. CHI (bottom path) pursues direct synthetic integration using hemispheric redundancy. Both aim to achieve a resilient Synthetic Consciousness Substrate (SCS) with the potential for shared technological advancements.}
    \label{fig:complementary-pathways}
  \end{figure}

  Purely synthetic substrates, as explored in CHI, offer the potential for inherent advantages over biological systems if, and only if, they can be rigorously proven capable of sustaining the necessary physical dynamics for consciousness: immunity to pathogens, tolerance for extreme environments (e.g., radiation, temperature), and engineered fault tolerance through redundancy. Synconetics remains deliberately agnostic to substrate taxonomy, recognising that biohybrid systems, blending biological resilience with synthetic fault tolerance (e.g., neural tissue interfaced with self-repairing nanowire meshes), may offer optimal transitional pathways. However, biology remains the only empirically validated substrate, a reality that ECP’s conservatism respects. Together, these approaches create a fail-safe continuum: synthetic substrate development proceeds without gambling existing consciousness, while biological augmentation extends survival until synthetic or hybrid options mature.

  \subsection{Engineering Resilience Through Substrate Design}
  \label{sec:engineering-resilience}

  Beyond the specific transition methodologies of ECP and CHI, a core objective of synconetics is the deliberate engineering of the target substrate (whether biohybrid or purely synthetic) for resilience far exceeding that of the original biological system. This enhanced robustness is not an incidental outcome but a primary design goal, achievable through specific architectural and operational advantages largely inaccessible to evolved biology. Key strategies enabling this resilience include decoupling conscious processes from bodily vulnerability, implementing distributed architectures for redundancy against localised damage, and facilitating high-resolution repair and maintenance.

  \begin{figure}[ht!]
    \centering
    \includegraphics[width=\textwidth]{figures/resilience-redundancy.png}
    \caption{\% TODO: Engineering Resilience via Distributed Architecture. (A) A centralised biological substrate is vulnerable to catastrophic failure from localised damage, terminating the process-world-line. (B) A synconetic substrate (SCS) designed with distributed redundancy can maintain processual continuity even if individual components fail.}
    \label{fig:resilience-redundancy}
  \end{figure}

  \textbf{Decoupling Agency from Embodiment:} Both methodologies enable the conscious substrate (SCS) to persist independently of its original biological body. Whether stored in a shielded facility (ECP’s ex vivo biohybrid brain) or distributed across synthetic nodes (CHI’s Artificial Hemispheric Complements), the SCS becomes decoupled from localised physical threats. Death of the biological body ceases to equate to termination of consciousness, akin to cloud computing surviving individual server failures.

  \textbf{Co-Location and Redundancy:} Synthetic or biohybrid substrates permit distributed architectures impossible in biology. Neural networks could span geographically isolated secure sites, with real-time synchronisation ensuring continuity. Partial destruction (e.g., 2\% loss from a localised attack) would not terminate the four-dimensional process-world-line, as critical dynamics persist across redundant nodes, mirroring RAID arrays in data storage.

  \textbf{Redundancy Through Resolution:} Modern computing achieves fault tolerance through atomic-scale repair (e.g., replacing individual transistors), yet biological brains lack analogous precision. High-bandwidth neural interfaces resolve this asymmetry, enabling synthetic repairs at resolutions matching the brain’s microcircuitry. This transforms the brain from a “black box” into an engineerable substrate, allowing repair and augmentation, such as enhancing memory density or computational speed while preserving continuity.

  \subsection{Sensory Abstraction as an Engineering Opportunity}
  \label{sec:sensory-abstraction}

  A further critical factor enabling the practical engineering of synconetic systems lies in the nature of sensory perception itself. Consciousness does not interface directly with raw physical reality but rather with processed, abstracted representations generated by biological sensory organs. Recognising and leveraging this inherent abstraction provides significant engineering opportunities for Synthetic Consciousness Interfacing (SCI).

  \textbf{Scale Separation of Sensory Input:} Biological consciousness already operates on abstracted sensory inputs: photons reduced to retinal signals, air vibrations to cochlear frequencies. SCI does not replicate raw physics; it only needs sufficient bandwidth to match the brain’s native compression. Virtual environments demonstrate this: triangular meshes evoke visceral fear of heights, and 44.1kHz audio fools humans into perceiving “live” music. By exploiting this abstraction, SCI systems can:

  \begin{itemize}
    \item Reduce engineering complexity (no need to simulate quantum fields),
    \item Enhance safety (virtual avatars avoid physical harm),
    \item Expand experiential range (perceiving infrared or ultrasonic ranges via synthetic transduction).
  \end{itemize}

  Existing sensory prosthetics (cochlear implants restoring hearing, retinal arrays granting vision, etc.) validate this approach. Their success hinges not on replicating biology exactly but on providing functionally equivalent inputs. Synconetics extends this principle, prioritising agency and coherence over perfect fidelity.

  \subsection{Stepwise Engineering of Substrate Continuity}
  \label{sec:pragmatic-path-forward}

  Given the potential for engineering substrate resilience and leveraging sensory abstraction, synconetics embraces a fundamentally pragmatic and stepwise approach to achieving substrate continuity. History demonstrates that transformative engineering often proceeds incrementally, frequently outpacing complete theoretical understanding; steam engines predated thermodynamics, and flight emerged before aerodynamic models. This tradition guides the focus on iterative manipulation of the known physical substrate (the brain, the CNS) rather than awaiting breakthroughs like a unified theory of consciousness and directly utilises rapid advancements across synergistic fields:

  \begin{itemize}
    \item \textbf{Neuroscience:} Mapping neural correlates and plasticity mechanisms.
    \item \textbf{Neuroengineering:} Developing high-resolution BCIs and stimulation protocols.
    \item \textbf{Materials Science:} Creating biocompatible interfaces and synthetic substrates.
  \end{itemize}

  Central to enabling such gradual interventions is the brain’s inherent plasticity. This capacity for functional reorganisation allows dynamic systems to adapt to structural change during interventions like ECP’s grafting or CHI’s hemispheric integration, aiming to maintain continuity much as occurs during recovery from certain neurological injuries.

  Nature’s evolutionary constraints yielded biological substrates optimised for survival, not engineered resilience. Just as aeronautics transcended avian flight (first mimicking feathers, then inventing turbines), synconetics aims to surpass biology’s limitations through deliberate engineering. Current efforts leveraging plasticity (ECP’s grafts, CHI’s AHCs) represent the necessary “propeller phase” of this trajectory: incremental improvements demonstrating feasibility. However, the end goal is the “jet engine” of consciousness substrates: systems built from potentially synthetic materials, incorporating fault-tolerant architectures, and operating based on physical principles extending beyond biology’s thermodynamic niche.

  \begin{figure}[ht!]
    \centering
    \includegraphics[width=\textwidth]{figures/engineering-progression.png}
    \caption{\% TODO: Stepwise Progression in Synconetic Engineering. Building on foundational sciences (Stage 1), initial approaches leverage biological mechanisms like plasticity (Stage 2: ‘Propeller Phase’). The trajectory aims towards deliberately engineered substrates with designed resilience surpassing biological limitations (Stage 3: ‘Jet Engine Phase’).}
    \label{fig:engineering-progression}
  \end{figure}

  Successfully navigating this stepwise path fundamentally reframes mortality. In this paradigm, death becomes both solvable and physically archaic: a contingent failure mode as avoidable as wooden biplanes in modern aviation. Preserving the four-dimensional process-world-line transitions from crisis management to routine maintenance, with engineered substrates offering near-invulnerability to traditional biological failure modes (trauma, ageing, disease). The ultimate engineering challenge then shifts from preventing collapse to optimising persistence across cosmological timescales, a feat demanding post-biological substrates and engineered transcendence.

  \subsection{Near-Term Impact and Iterative Progress}
  \label{sec:near-term-impact}

  Synconetics’ feasibility is amplified by its alignment with near-term technological trajectories. Advances in neural prosthetics, organoid development, and biocompatible materials (already funded for medical applications) directly serve substrate stabilisation and repair. Early milestones, such as restoring motor function via BCIs or reversing age-related neural decline, offer tangible societal benefits, ensuring continued investment even before full continuity is achieved.

  The engineering challenge, while immense, decomposes into tractable sub-problems:
  \begin{itemize}
    \item Validating synthetic substrates (CHI’s focus),
    \item Ensuring continuity during the transition (ECP’s gradualism),
    \item Optimising sensory interfaces (shared SCI development).
  \end{itemize}

  This modularity allows parallel progress. By prioritising incremental, testable advancements over speculative leaps, synconetics transforms the pursuit of an indefinite lifespan into an engineering roadmap.

  % ! ================================
  % ! # MARK: Roadmap and Implications
  % ! ================================

  \section{Roadmap and Implications}
  \label{sec:roadmap}

  Synconetics’ viability rests on a pragmatic roadmap acknowledging both profound challenges and near-term opportunities. The dual methodologies presented (ECP pursuing gradual biohybrid transition and CHI exploring synthetic substrate integration) represent complementary strategies within this framework (\autoref{sec:methods}). ECP leverages current biological understanding, offering potential near-term resilience gains. At the same time, CHI directly confronts the engineering of synthetic substrates, a path promising greater long-term robustness but facing significant hurdles in validation. This parallel pursuit provides strategic hedging: advances in one inform the other (e.g., BMI development, substrate validation criteria), mitigating risks associated with the inherent uncertainties in substrate science. Crucially, the potential success of either approach within a foreseeable timeframe, possibly the next two decades, compels us to confront the profound implications not as distant philosophical exercises but as urgent practical concerns demanding immediate, serious consideration.

  The potential emergence of individuals persisting via engineered substrates (i.e., synconetic entities) fundamentally challenges our existing societal, legal, and ethical frameworks, raising critical questions across every domain. Current legal systems, predicated entirely on biological personhood, are unprepared. How would legal identity, rights (such as property ownership, voting, substrate autonomy), and responsibilities (taxation, liability) be defined for an entity potentially lacking a conventional body? Does personhood track the continuous process-world-line irrespective of the substrate, as synconetics principles suggest? Establishing internationally recognised standards and dynamic legal frameworks capable of adapting to substrate transitions and potential capability augmentations represents a monumental undertaking, essential to prevent inequality and novel forms of exploitation.

  Economically, the advent of potentially vastly long-lived individuals promises radical disruption. How would Synconetic entities participate in labour markets increasingly shaped by AI automation? What forms of value could they create, perhaps leveraging enhanced cognitive capabilities or unique perspectives gained from virtual environments? Does their existence necessitate fundamental shifts in economic models, potentially accelerating discussions around systems like Universal Basic Income? Furthermore, the significant, ongoing resource demands for sustaining engineered substrates (energy, computation, physical security, and specialised maintenance) raise critical questions of allocation and equity. What economic models (e.g., subscription, public utility, private ownership) could govern access and upkeep without creating unprecedented societal stratification based on the ability to afford continued existence?

  Infrastructurally and logistically, supporting a population of Synconetic entities involves daunting engineering challenges. What constitutes robust, secure physical and digital infrastructure for consciousness-supporting substrates? Where would they be housed, and what levels of physical security, redundancy against failure or catastrophe, and resilience against malicious attack are achievable, sustainable, and ethically mandated? Do centralised hosting models create unacceptable single points of failure and control, and are distributed architectures practically feasible while maintaining real-time synchronisation for continuity? Critically, what safeguards ensure provider viability and prevent vendor lock-in? Without clear standards guaranteeing substrate portability and continuity during provider transitions, individuals face extreme vulnerability. Can the significant energy and computational load, particularly for rich virtual interactions facilitated by SCI, be sustained globally?

  Perhaps most profoundly, the successful realisation of synconetics challenges fundamental societal notions of life, death, identity, and community. How will society perceive Synconetic entities: as extensions of their former biological selves, as ‘post-biological’, or as something entirely new? How must relationships, family structures, inheritance laws, and social support systems adapt? What are the psychological implications for individuals undergoing transition and potentially existing indefinitely, possibly within non-biological environments? How can existential meaning and well-being be maintained under such radically different conditions? Ensuring equitable access, mitigating potential coercion (e.g., societal pressure to transition), and defining death for a synconetic entity, including ethically sound protocols for managing substrate failure, irreversible cognitive decline, or respecting a voluntary wish to cease existence, represent entirely uncharted territories demanding sensitive, cross-disciplinary deliberation now.

  \begin{figure}[ht!]
    \centering
    \includegraphics[width=\textwidth]{figures/readiness-gap.png}
    \caption{\% TODO: Potential Readiness Gap in Synconetics Development. Technical feasibility (Curve 1) may advance more rapidly than the necessary societal, legal, and ethical frameworks (Curve 2). Proactive, parallel development of these frameworks is crucial to mitigate risks associated with this gap and ensure responsible innovation.}
    \label{fig:readiness-gap}
  \end{figure}

  The near-term feasibility targeted by synconetics transforms these profound questions from speculative fiction into urgent matters for contemporary policy, ethics, and engineering design. Unlike paradigms focused on indefinite futures, synconetics’ potential impact within the professional lifetime of researchers working today underscores the imperative for proactive engagement. Addressing the legal, economic, infrastructural, and ethical dimensions cannot be postponed; it must occur in parallel with technical research and development. Proposing mechanisms like an International Continuity Commission, open-source substrate architectures, or mandatory continuity impact assessments serve to illustrate the kind of proactive governance needed. This foresight is essential to mitigate the risks of societal disruption, inequality, and catastrophic failure, ensuring that the pursuit of engineered continuity aligns with broadly shared human values and responsible innovation, a core tenet of the synconetics approach.

  % ! ==================
  % ! # MARK: Conclusion
  % ! ==================

  \section{Conclusion}
  \label{sec:conclusion}

  Synconetics reframes biological death as a tractable engineering problem, solvable through the rigorous pursuit of verifiable processual continuity. This essay has laid the groundwork; translating these principles into reality demands immediate, concerted action. The methodologies presented demonstrate feasibility, but accelerating progress requires a deeply collaborative, transdisciplinary effort that extends beyond our initial working group.

  We, therefore, issue a direct call to action. For researchers seeking to advance the fundamental science and engineering, opportunities exist for postgraduate work exploring consciousness mechanisms and BMI integration (e.g., with Prof. Watanabe’s group at the University of Tokyo) or for developing other unexplored continuity-preserving strategies. For engineers and scientists focused on tangible applications, organisations like Eightsix Science are actively seeking technical collaborators, funding, and grant support to translate synconetics principles into practice via approaches like ECP.

  Beyond these specific avenues, we urge innovators to foster diversity by launching new research projects or companies exploring alternative SCM pathways. We call upon funders to recognise the near-term potential and support ventures grounded in empirical validation and continuity. We invite ethicists, legal scholars, policymakers, and social scientists to engage proactively with the profound implications outlined, helping to shape responsible governance. We encourage critical dialogue: connect with the authors, challenge assumptions, join nascent community discussions (e.g., via our Discord server), and help refine this framework. Disseminate this work, engage peers, and contribute to future knowledge-building efforts like a potential planned book on these topics.

  Synconetics is not a distant theoretical exercise; it is an engineering imperative demanding focus and collaboration today. By concentrating on verifiable physical continuity and buildable systems, we can move beyond speculation towards tangible progress. We welcome all who share this commitment to join us in building this critical field.

  % ! ========================
  % ! # MARK: References, etc.
  % ! ========================

  \pagebreak
  \bibliographystyle{../../templates/custom-apa}
  \bibliography{references/bibliography}
  \nocite{*}

\end{sloppypar}
\end{document}